% ========================================================
% This document is a customizable CV/Resume template built using LaTeX.
% The template is designed for easy customization and clear structure.
%
% Author: Matthew DeVerna (www.matthewdeverna.com)
% Originally Published Date: 2024
% Design: Inspired by Hause Lin's CV (hauselin.com), which has changed a good deal since this was created.
% 
% Change Log:
% -----------------
% - 2025-08-19: Include reverse numbering for bibliography sections
% - 2026-08-29: Fix \cite printing the mirror image of the reversed list
%               labels (e.g. citing J6 printed "J2" in a 6-item list)
% 
% Project Overview:
% -----------------
% This LaTeX document is designed to help you create a professional CV or 
% resume with ease. It uses as little fancy LaTex functionality or custom functions as possible to maximize its longterm durability and flexibility.
% The document is structured into multiple sections, each loaded from 
% separate subfiles for modularity and ease of maintenance. 
%
% Key Features:
% -------------
% - Customizable sections: Education, Research Experience, Awards, Publications, etc.
% - Bookmarks in the PDF for easy navigation
% - Styled bibliography with BibLaTeX
% - Hyperlinked email and website
% - FontAwesome icons for additional styling
%
% Getting Started:
% ----------------
% 1. Customize your personal information by modifying the \mytitle command.
% 2. Add your content to the respective subfiles (e.g., education.tex, exp_research.tex).
% 3. Update the bibliography file (ref.bib) with your publications and categorize them with keywords.
%
% Important Notes:
% ----------------
% - This main file includes the overall structure and settings. 
% - Each section has its own detailed instructions for further customization.
%
% ========================================================

\documentclass[11pt]{article} % Choose the document class and font size
\usepackage[margin=1in]{geometry} % Page layout settings

% Set the citation style
\usepackage[
    backend=biber,      % Specifies the backend to be used by BibLaTeX for processing the bibliography. 'biber' is the default backend.
    maxnames=20,        % Limits the maximum number of author names to display before abbreviating with "et al."
    style=nature,       % Sets the citation style to 'nature,' which is commonly used in scientific papers.
    sorting=ydnt,       % Specifies the sorting order of entries in the bibliography:
                        % y - year (descending)
                        % d - descending order
                        % n - name
                        % t - title
    defernumbers=true,  % Delays the assignment of citation numbers until the end of the document, allowing for the correct order of citations within each bibliography section.
]{biblatex}
\addbibresource{ref.bib} % Adds the bibliography resource file 'ref.bib' containing all the references.

% -- Reverse-numbered bibliography (3, 2, 1) environment START ---
\usepackage{etaremune}
\defbibenvironment{revbib}
  {%
    \csname begin\endcsname{etaremune}%
    % Align indent with biblatex’s numeric label width
    \setlength{\leftmargin}{\dimexpr\labelnumberwidth+\labelsep\relax}%
    % Prefix from \newrefcontext[labelprefix=...] + descending number (no parentheses)
    \renewcommand*{\labelenumi}{\printfield{labelprefix}\arabic{enumi}}%
  }
  {%
    \csname end\endcsname{etaremune}%
  }
  {\item}
% -- Reverse-numbered bibliography environment END ---
% NOTE: If you want the publications list in each section to start from 1 and count up (1, 2, 3), then you can simply remove the `env=revbib` line in all of the \printbibliography lines.

% -- Keep \cite in sync with the reversed numbers START ---
% WHY THIS IS NEEDED:
% `etaremune` numbers the printed list downwards (N, ..., 2, 1), but BibLaTeX's
% own `labelnumber` field still counts upwards (1, 2, ..., N) in that same
% order. Because \cite prints `labelnumber`, an uncorrected \cite emits the
% MIRROR IMAGE of the label you can see in the list: in a 6-item list, citing
% the item labelled J6 would print "J2", and so on.
%
% The fix: record how many entries each labelprefix list holds (written to the
% .aux file while the list is typeset), then print (N + 1 - labelnumber)
% wherever \cite asks for a number. The visible list labels are produced by
% etaremune and are left untouched.
%
% NOTE: because the counts travel through the .aux file, a fresh build needs
% two LaTeX passes before citations settle. `latexmk` handles this for you.
% If you removed `env=revbib` above (i.e. you are NOT reverse numbering), you
% should delete this block too.
\makeatletter
% Storage for the per-prefix totals, replayed from the .aux file each run.
\providecommand*{\revbibsetcount}[2]{}
\renewcommand*{\revbibsetcount}[2]{%
  \expandafter\xdef\csname revbib@count@#1\endcsname{#2}%
}
% Entries are typeset 1..N, so the last value written for a prefix equals N.
\AtEveryBibitem{%
  \iffieldundef{labelprefix}
    {}
    {\savefield{labelprefix}{\revbib@pfx}%
     \savefield{labelnumber}{\revbib@num}%
     \protected@write\@auxout{}{%
       \string\revbibsetcount{\revbib@pfx}{\revbib@num}}}%
}
% Turn an ascending number into the descending one the reader actually sees.
\newcommand*{\revbib@reverse}[2]{%
  \ifcsname revbib@count@#1\endcsname
    \number\numexpr\csname revbib@count@#1\endcsname+1-#2\relax
  \else
    #2% no count recorded yet (first pass) -- fall back to the raw number
  \fi}
\DeclareFieldFormat{labelnumber}{\revbib@reverse{\thefield{labelprefix}}{#1}}
\makeatother
% -- Keep \cite in sync with the reversed numbers END ---

% Allows columns that stretch across pages
\usepackage{longtable}

% Table functionality and beautification (not strictly needed)
\usepackage{bookmark}

% Use icons, if you want.
% All available icons: http://mirrors.ibiblio.org/CTAN/fonts/fontawesome5/doc/fontawesome5.pdf
\usepackage{fontawesome}

% Allows font justification control (needed for clean pub-list formatting)
\usepackage{ragged2e}

% For underlining with line breaks
\usepackage{soul} 

% All fonts: https://tug.org/FontCatalogue/
\usepackage{kpfonts} % More professional font
% \usepackage[default]{sourcecodepro} % Code-like font
\usepackage[T1]{fontenc}

% Control hyperlinks and colors
% CUSTOM COLORS INCLUDED DIRECTLY AFTER \begin{document}
\usepackage{xcolor}
\usepackage{hyperref}
\hypersetup{
    colorlinks=true,        % Enable colored links
    breaklinks=true,        % Allow links to break across lines
    linkcolor=cornflowerblue,    % Color of internal links
    urlcolor=cornflowerblue,     % Color of URL links
    anchorcolor=cornflowerblue,  % Color of anchors
    citecolor=cornflowerblue,    % Color of citations
    pdftitle={Your Title},    % Title of the PDF
    pdfauthor={Your Name}, % Author of the PDF
    bookmarksopen=true,      % Open bookmarks panel at start
}

%%% CONVENIENCE FUNCTIONS GO HERE %%%
%%% ----------------------------- %%%
\newcommand{\mytitle}[4]{
  \begin{center}
    \Large\textbf{#1}\normalsize \\ % Name in large bold font
    \href{mailto:#2}{#2} \\ % Email with mailto: link
    \href{https://#3}{#3} \\ % Website with link
    #4 % Address
  \end{center}
}
%%% ----------------------------- %%%


\begin{document}
\include{colors.tex} % Load custom colors from colors file
\mytitle{Your Name}{person@gmail.com}{yourwebsite.com}{Your Address\\Goes here\\This field can be excluded} % Insert your custom title


% Ensure right side margin is not surpassed by bibliography and the right margin is aligned throughout
\RaggedRight


% These \pdfbookmark lines create bookmarks in the exported PDF document that display in the left pane.
% Value in [] sets the indentation level of the bookmark
\pdfbookmark[1]{Education}{}
\section*{Education}
% Add your educational background here!

% NOTE: If you want to remove the "Expected" footnote, you will want to remove:
% - Directly below: \renewcommand, \setcounter
% - In the table: \footnotemark in the left column
% - After the table: \footnotetext, \renewcommand, \setcounter

% Different numbers in "\setcounter{footnote}{0}" use different symbols
\renewcommand{\thefootnote}{\fnsymbol{footnote}}
\setcounter{footnote}{0}

\begin{longtable}[l]{@{}p{.1\textwidth} p{\dimexpr.9\textwidth-2\tabcolsep\relax}@{}}
    % Use custom symbol footnote for "expected"
    2020--25\footnotemark & \textbf{Ph.D.}, Your Field, Your University \\

    2018-20 & \textbf{M.A.}, Field, University \\

    2008-12 & \textbf{B.Sc.} Field, University
% \end{tabularx}
\end{longtable}


% Add text for the custom footnote
\footnotetext[1]{Expected.}

% Restore the default footnote numbering
\renewcommand{\thefootnote}{\arabic{footnote}}
\setcounter{footnote}{1}


\pdfbookmark[1]{Research Experience}{exp_research}
\section*{Research Experience}
\label{exp_research}
% List your research experience here

\begin{longtable}[l]{@{}p{.1\textwidth} p{\dimexpr.9\textwidth-2\tabcolsep\relax}@{}}

    2020-- & Research Assistant, \href{https://example.com/}{Cool Lab}, Your University, (Advisor: \href{https://example.com/}{Dr. Thinks-a-lot}) \\

    2020-- & Research Assistant, \href{https://example.com/}{Fancy Science Lab}, Your University, (Advisor: \href{https://example.com/}{Dr. Math Wiz}) \\

\end{longtable}

\pdfbookmark[1]{Awards \& Honors}{awards}
\section*{Awards \& Honors}
\label{awards}
% Include awards, honors, grant funding, etc. here

\begin{longtable}[l]{@{}p{.1\textwidth} p{\dimexpr.9\textwidth-2\tabcolsep\relax}@{}}
    % Use custom symbol footnote for "expected"
    2024 & \href{https://www.nsf.gov/}{Fancy NSF Grant} (\$1,000,000; Role: PI) \\
   
    2023 & Your University Outstanding Service Award (\$XXX) \\
       
    2020-- & Research Fellow at the \href{https://example.com/}{Your Research Lab or University} \\

    % Add more things here!
   
\end{longtable}



\pdfbookmark[1]{Publications}{pubs}
\section*{Publications}
\label{pubs}

% Add equal contribution dagger
\vspace{-.75em}
\small
\faGoogle~\href{https://scholar.google.com/}{Google Scholar}\\
$\dagger \rightarrow$ Equal contribution
\normalsize


\pdfbookmark[2]{Journal Articles}{journal-article}
\subsection*{Journal Articles}
\label{journal-article}
\newrefcontext[labelprefix=J] % Will prefix bibliography numbers with this letter
% Ensures publications which are not cited in the document are included in the above sections
\nocite{*} % Ensures uncited items are included
\printbibliography[
    type=article, % Only include @article ref.bib items
    heading=none, % Do not include header. Gives us more control.
    resetnumbers=true, % Start item counter from zero
    keyword=J, % Include items in ref.bib with keyword={J}
    env=revbib
]

\pdfbookmark[2]{Peer-reviewed Conference Proceedings}{conferences}
\subsection*{Peer-reviewed Conference Proceedings}
\label{conferences}
\newrefcontext[labelprefix=C]
\printbibliography[type=inproceedings,heading=none,resetnumbers=true,keyword=C, env=revbib]

\pdfbookmark[2]{Working papers}{working-papers}
\subsection*{Working papers}
\label{working-papers}
\newrefcontext[labelprefix=W]
\printbibliography[type=misc,heading=none,resetnumbers=true,keyword=W, env=revbib]
    

\pdfbookmark[1]{Tools \& Software}{tools}
\section*{Tools \& Software}
\label{tools}
\input{tools.tex}


\pdfbookmark[1]{Presentations}{presentations}
\section*{Presentations}
\label{presentations}

% Include any additional details here
% \vspace{-.75em}
% \small
% $\dagger \rightarrow$ Equal contribution
% \normalsize

\pdfbookmark[2]{Talks}{talks}
\subsection*{Talks}
\label{talks}
\newrefcontext[labelprefix=T]
\printbibliography[type=misc,heading=none,resetnumbers=true,keyword=T, env=revbib]

\pdfbookmark[2]{Posters}{posters}
\subsection*{Posters}
\label{posters}
\newrefcontext[labelprefix=P]        
\printbibliography[type=misc,heading=none,resetnumbers=true,keyword=P, env=revbib]

\pdfbookmark[2]{Demonstrations \& Tutorials}{demos}
\subsection*{Demonstrations \& Tutorials}
\label{demos}
\newrefcontext[labelprefix=D] \printbibliography[type=misc,heading=none,resetnumbers=true,keyword=D, env=revbib]


\pdfbookmark[1]{Selected Media Coverage}{media}
\section*{Selected Media Coverage}
\label{media}
%  Include media coverage here
% 
\begin{longtable}[l]{@{}p{.1\textwidth} p{\dimexpr.9\textwidth-2\tabcolsep\relax}@{}}

    2024 & The New York Times, \href{https://nyt.com}{Scientists figure out the meaning of life.}~\cite{art1} (Note how the regular citation functions work throughout the document!) \\
    
    2024 & AP News, \href{https://apnews.com/}{What even is science?}~\cite{art2} (Point to any project in your ref.bib file with the standard `cite' function) \\

    2025 & Reuters, \href{https://reuters.com}{The newest paper makes headlines.}~\cite{art3} (This one cites the \emph{top} journal entry, so it prints J3 -- matching the number shown in the list above) \\

    
\end{longtable}

 % Input lines load the material from the subdocuments


\pdfbookmark[1]{Teaching}{teaching}
\section*{Teaching}
\label{teaching}

\subsection*{University1 Name}
% 
\begin{longtable}[l]{@{}p{.1\textwidth} p{\dimexpr.9\textwidth-2\tabcolsep\relax}@{}}

    2024 & Assistant Instructor, Course Name (CODE 606) \\
    
    2024 & Guest lecturer, Science Theory and Practice (CODE 360) \\

\end{longtable}


\subsection*{University2 Name}
% 
\begin{longtable}[l]{@{}p{.1\textwidth} p{\dimexpr.9\textwidth-2\tabcolsep\relax}@{}}

    2024 & Assistant Instructor, Course Name (CODE 123) \\
    
    2024 & Guest lecturer, Science Theory and Practice (CODE 321) \\

\end{longtable}


\pdfbookmark[1]{Academic Advising}{advising}
\section*{Academic Advising}
\label{advising}
% Including academic advisement history here


% Remove \subsection{} lines and multiple tables if you only need one section!
\subsection*{Graduate}
\begin{longtable}[l]{@{}p{.1\textwidth} p{\dimexpr.9\textwidth-2\tabcolsep\relax}@{}}

    2023 & Name, University, Additional details \\

\end{longtable}

\subsection*{Undergraduate}
\begin{longtable}[l]{@{}p{.1\textwidth} p{\dimexpr.9\textwidth-2\tabcolsep\relax}@{}}

    2023 & Name, University, Additional details \\

\end{longtable}


\pdfbookmark[1]{Academic Service}{service}
\section*{Academic Service}
\label{service}


\subsection*{Guest Editor}
% 
\begin{longtable}[l]{@{}p{.1\textwidth} p{\dimexpr.9\textwidth-2\tabcolsep\relax}@{}}


    2023 & \href{https://example.com}{Link to collection} \\

% \end{tabularx}
\end{longtable}


\subsection*{Journal Reviewer}
% 
\begin{longtable}[l]{@{}p{.1\textwidth} p{\dimexpr.9\textwidth-2\tabcolsep\relax}@{}}

    2023 & Nature\\
    2023 & Science \\
    
% \end{tabularx}
\end{longtable}





\pdfbookmark[1]{Other Experience}{exp_other}
\section*{Other Experience}
\label{exp_other}
% Add other experience you don't know where to classify here

\begin{longtable}[l]{@{}p{.1\textwidth} p{\dimexpr.9\textwidth-2\tabcolsep\relax}@{}}

    2018 & Past Non-academic Job, \href{https://example.com/}{Company Name}, New York City, New York (Additional detail) \\

\end{longtable}


% Pretty ending with the date last updated
\centering
\rule{0.25\linewidth}{0.4pt}\\
\medskip
Last updated: \today

\end{document}
